\chapter{Valores, completitud y bloques de la jerarquía espectral}
\label{chap:torres}

El operador bicapa, sus momentos y el semigrupo
\(\mathfrak S=\langle90,120\rangle\) se definen una sola vez en el
\cref{sec:jerarquia-espectral-infinita-rev7}. Este capítulo publica valores,
prueba la completitud reconstructiva y separa los bloques globales. Los
mismos enteros \(90,120\) aparecen también como cardinalidades en la
incidencia hexada--octada, pero esa coincidencia no transporta la incidencia
\(\iota_{6,8}\) a la carta angular ni convierte sus configuraciones en
momentos de masa.

\section{Primeras evaluaciones de la jerarquía}

Los índices
\[
 90,120,180,210,240,270,360,420
\]
son una selección de primeras alturas de \(\mathfrak S\), pues
\[
 180=2\cdot90,\quad210=90+120,\quad240=2\cdot120,
\]
\[
 270=3\cdot90,\quad360=3\cdot120,
 \quad420=2(90+120).
\]
Los momentos \(s_n=q_-^n-q_+^n\) y la interfaz generadora
\(\mathcal K_{\mathrm{ang}}^{90,120}=(s_{90},s_{120})\) se usan aquí con la
convención ya fijada; no se redefinen. La totalidad de los valores queda
determinada por \(A\), \(\Cstar\) y la recurrencia cuadrática del lector.

\begin{table}[htbp]
\centering
\caption{Momentos orientados de la jerarquía global.}
\label{tab:torres}
\begin{tabular}{@{}rr@{}}
\toprule
\(n\)&\(s_n\)\\
\midrule
90  & \(2.9516943928982796\times10^{-4}\)\\
120 & \(1.9683511250294992\times10^{-5}\)\\
180 & \(8.7345871869508335\times10^{-8}\)\\
210 & \(5.8181313057291163\times10^{-9}\)\\
240 & \(3.8754674969305336\times10^{-10}\)\\
270 & \(2.5814553607509555\times10^{-11}\)\\
360 & \(7.6293257871406076\times10^{-15}\)\\
420 & \(3.3850663628348497\times10^{-17}\)\\
\bottomrule
\end{tabular}
\end{table}

\begin{proposition}[Carácter global de los momentos]
Las cantidades \(s_n\) de la \cref{tab:torres} son invariantes de la carta
angular fijada. Una especie aporta coordenadas enteras que ponderan esos
invariantes, pero no modifica \(q_\pm\) ni introduce una torre propia.
\end{proposition}
\begin{proof}
Cada \(s_n\) es una función de \(q_-\) y \(q_+\). El teorema de inversión de
canales demuestra que la aplicación conjunta
\((A,\Cstar)\mapsto(q_+,q_-)\) es inyectiva. No aparece ninguna etiqueta de
especie en estas definiciones.
\end{proof}

\begin{theorem}[Completitud reconstructiva de la jerarquía]
\label{thm:completitud-reconstructiva-torres}
Sea \(n_1<n_2<\cdots\) la enumeración creciente sin repeticiones de
\(\mathfrak S\). Para
\[
 a_k:=s_{n_k}=q_-^{n_k}-q_+^{n_k}>0,
\]
toda razón positiva \(\rho\in\RR_{>0}\) admite una expansión absolutamente
convergente en la jerarquía:
\[
 \boxed{\rho=
 \exp\!\left(\sum_{k\geq1}\nu_k a_k\right)}
 \qquad\text{para alguna sucesión }(\nu_k)_{k\geq1}\subset\ZZ.
\]
\end{theorem}

\begin{proof}
Como \(0<q_+<q_-<1\), se cumple
\[
 0<a_k<q_-^{n_k}.
\]
Los \(n_k\) son múltiplos de \(30\) no menores que \(90\); por tanto,
\[
 \sum_{k\geq1}a_k
 \leq\sum_{m\geq3}q_-^{30m}
 =\frac{q_-^{90}}{1-q_-^{30}}<\infty,
 \qquad a_k\longrightarrow0.
\]
Fijemos \(r_0=\log\rho\) y una regla de desempate para el entero más próximo.
Definimos recursivamente
\[
 \nu_k=\operatorname{nint}\!\left(\frac{r_{k-1}}{a_k}\right),
 \qquad
 r_k=r_{k-1}-\nu_k a_k.
\]
La propiedad del entero más próximo da
\[
 |r_k|\leq\frac{a_k}{2}\longrightarrow0.
\]
La identidad telescópica
\[
 r_0=\sum_{j=1}^{N}\nu_j a_j+r_N
\]
implica \(r_0=\sum_{j\geq1}\nu_j a_j\). La convergencia es absoluta, pues
\[
 \begin{aligned}
 \sum_{k\geq1}|\nu_k|a_k
 &=\sum_{k\geq1}|r_{k-1}-r_k|\\
 &\leq |r_0|+\frac{a_1}{2}
 +\frac12\sum_{k\geq2}(a_{k-1}+a_k)<\infty.
 \end{aligned}
\]
Al exponenciar se obtiene la afirmación.
\end{proof}

\begin{remark}[Alcance y control negativo]
El teorema es una sobreyectividad reconstructiva del lenguaje de torres sobre
\(\RR_{>0}\): prueba resolución del codominio, no selección física. El
algoritmo recibe \(r_0=\log\rho\); si \(\rho\) es una magnitud objetivo, sus
coeficientes son, por ello, una reconstrucción dependiente del objetivo. La
procedencia prospectiva requiere otra flecha,
\[
 \text{especie}\dashrightarrow\text{extensión}
 \longrightarrow\text{ruta}\longrightarrow\text{registro}
 \longrightarrow\text{firma},
\]
y no se deduce de esta expansión. La conclusión fallaría si los \(a_k\) no
tendieran a cero, si la cota del residuo dejara de cumplirse o si la serie de
errores no fuera sumable.
\end{remark}

\section{Bloques globales de la acción de masa}

La acción de masa utilizará la terna
\[
 \Dfour,\qquad s_{120},\qquad \zCorr=135\Dfour^2.
\]
Sus valores se calculan una sola vez. Las coordenadas
\(k,\nu_{120},\nu_{270}\) de una firma especifican cuántas veces interviene
cada bloque en el exponente; no cambian su valor.

Conviene distinguir dos objetos que comparten el índice \(270\):
\[
 \boxed{\zCorr=135\Dfour^2
 =1.66922699663643\times10^{-6}}
\]
y
\[
 \boxed{\sSpec=q_-^{270}-q_+^{270}
 =2.58145536075096\times10^{-11}.}
\]
El primero es un invariante cuadrático construido a partir de \(\Dfour\); el
segundo es un momento espectral del semigrupo. La coincidencia del subíndice no
autoriza su identificación.

\section{Combinación de orden superior}

La combinación global
\[
 \delta\Dfour^{\rm tower}
 =\frac{s_{180}}{24}-\frac{s_{240}}{54}
 -24s_{360}+\frac76s_{420}
\]
evalúa a
\[
 \delta\Dfour^{\rm tower}
 =3.632051471908759123889070967\times10^{-9}.
\]
La coordenada de comparación utilizada en la selección de los
coeficientes se define como el dato escalar
\[
 \delta\Dfour
 :=3.632051471908972784661641720\times10^{-9}.
\]
Por consiguiente, el residuo de la combinación respecto de esa coordenada es
\[
 \delta\Dfour^{\rm tower}-\delta\Dfour
 =-2.1366077257075\times10^{-22}.
\]
La cantidad pequeña es, pues, una diferencia entre dos objetos y no el valor
de la combinación. Todos los sumandos pertenecen a la misma jerarquía y sus
coeficientes son globales. Como la coordenada de comparación intervino en la
selección de los coeficientes, esta igualdad es reconstructiva. La
combinación no se añade de nuevo a la fórmula electrónica que ya emplea
\(\Dfour\), pues esa sustitución duplicaría la misma corrección.
